\chapter{Callbacks, Reentrancy, and Threads}
\label{ch:ffi-callbacks}

A callback reverses the direction of control before the original foreign operation is finished.
The native library may call managed code while holding a lock or while its own object is temporarily inconsistent.
That managed code may call the library again, close an object, or start another operation.
Correct register transfers alone cannot make this control flow safe.

\section{Pair executable entry with retained context}
\label{sec:ch73-context}
A \term{callback registration} associates an executable entry point with the state required to invoke a particular callable.
A C interface taking both a function pointer and a context pointer can use one static trampoline for many callback objects.
The context identifies the runtime callable and any error or lifetime state.
An interface accepting only a function pointer may instead need a generated trampoline, such as a libffi closure, whose entry recovers the context~\cite{libffi-2026-manual}.

The running visitor accepts a callback, context, and count.
For indices $0$ through $n-1$, the callback writes a signed 32-bit value and returns a status.
The visitor accumulates values into a signed 64-bit sum and stops at the first failure.
Its bound $0\le n\le10{,}000$ guarantees that summing any allowed callback results fits that accumulator.
Figure~\ref{fig:ch73-roundtrip} shows the nested control path.

\begin{figure}[htbp]
\centering
\begin{tikzpicture}[node distance=12mm]
\node[basic block] (managed) {managed visitor call};
\node[basic block,below=of managed] (native) {native loop};
\node[basic block,below=of native] (trampoline) {protected callback adapter};
\node[basic block,below=of trampoline] (body) {managed callable};
\draw[cfg edge] (managed)--(native); \draw[cfg edge] (native)--(trampoline);
\draw[cfg edge] (trampoline)--(body);
\path (body.east) ++(3,0) coordinate (outer);
\draw[cfg edge] (body.east)--(outer) |- (native.east);
\node[right,font=\small] at (outer |- trampoline) {may reenter};
\end{tikzpicture}
\caption{A callback can start another foreign operation before the first one returns.}
\label{fig:ch73-roundtrip}
\end{figure}

For the callback $f(i)=i+2$ and $n=5$, the values are $2,3,4,5,6$ and the sum is $20$.
If the callback fails at index $2$, the visitor has observed two successful values but does not publish their partial sum.
It returns a failure with the caller's output unchanged.
Table~\ref{tab:ch73-trace} records both paths.

\begin{table}[htbp]
\caption{The visitor publishes only a complete successful aggregate.}
\label{tab:ch73-trace}
\centering
\begin{tabular}{rrrrl}
\toprule
Index & Value & Success sum & Failure-path sum & Failure-path action \\
\midrule
0 & 2 & 2 & 2 & continue \\
1 & 3 & 5 & 5 & continue \\
2 & 4 & 9 & 5 & stop; preserve output \\
3 & 5 & 14 & -- & not called \\
4 & 6 & 20 & -- & not called \\
\bottomrule
\end{tabular}
\end{table}

\section{Preserve reentrant invariants}
\label{sec:ch73-reentrant}
The case study uses no global callback slot and holds no library mutex while calling managed code.
Each invocation has its own context on the calling stack.
The callable remains rooted by the original language arguments throughout the synchronous call.
Nested visits therefore have separate counters, sums, and error state.

This design supports reentrancy because the library's visible state is already consistent at the callback point.
A production container visitor must choose a policy for mutation during iteration: prohibit it, detect a version change, snapshot the iteration set, or define stable iterator semantics.
Releasing a mutex before calling user code avoids one deadlock pattern, but does not by itself preserve the container's iteration invariants.
Reacquiring it afterward may reveal that the object has changed or closed.

Algorithm~\ref{alg:ch73-callback} states the synchronous callback protocol.
It costs $O(n)$ invocations and $O(1)$ native bookkeeping, excluding callback cost and recursion initiated by user code.
The count bound limits this loop, not arbitrary computation inside the callback.

\begin{algorithm}[htbp]
\caption{Visit indices with a synchronous status-returning callback.}
\label{alg:ch73-callback}
\begin{pseudocode}
\Require bounded count $n$, rooted callable, attached runtime thread, writable output
\Ensure full sum on success; output unchanged on failure
\State validate $0\le n\le10{,}000$; $s\gets0$
\For{$i=0$ to $n-1$}
  \State invoke callback through its runtime error boundary
  \If{callback failed or result is outside int32 range}
    \State return callback failure without publishing $s$
  \EndIf
  \State $s\gets s+\text{callback result}$
\EndFor
\State publish $s$; return success
\end{pseudocode}
\end{algorithm}

\section{Make registration teardown wait for users}
\label{sec:ch73-quiescence}
Asynchronous registration requires stronger ownership than a synchronous stack context.
The library can retain the callable and invoke it after registration returns.
Unregistering must establish that no new invocation can start and that every previously admitted invocation has completed before freeing the context or trampoline.
\term{Callback quiescence} is this absence of both future admission and active users.

Figure~\ref{fig:ch73-stop} shows the legal retirement path.
An atomic count alone is insufficient if a thread can read a raw context pointer before another thread frees it and only then try to increment the count.
Admission must occur through a registry or other stable gate whose own lifetime outlasts every attempted admission.
Algorithm~\ref{alg:ch73-retire} assumes such a gate and serializes admission and stopping under its mutex.

\begin{figure}[htbp]
\centering
\begin{tikzpicture}[node distance=12mm]
\node[basic block] (open) {accepting; active count may grow};
\node[basic block,below=of open] (stop) {stopped; active users finish};
\node[basic block,below=of stop] (free) {active count zero; reclaim};
\draw[cfg edge] (open)--(stop); \draw[cfg edge] (stop)--(free);
\end{tikzpicture}
\caption{Stopping admission precedes draining users and reclaiming callback state.}
\label{fig:ch73-stop}
\end{figure}

\begin{algorithm}[htbp]
\caption{Retire a registration through a stable serialized admission gate.}
\label{alg:ch73-retire}
\begin{pseudocode}
\Require stable gate shared by every admission; caller is not its own active callback
\Ensure no callback can use the reclaimed context
\State under gate mutex: mark registration as not accepting
\State remove it from future library dispatch under the library's unregister contract
\State wait for active count to become zero, releasing the mutex while waiting
\State detach context and trampoline from the registry
\State release runtime root and executable closure storage through their owners
\end{pseudocode}
\end{algorithm}

Bookkeeping is constant per admission and completion; retirement latency includes the longest outstanding callback and has no finite bound if a callback never returns.
Calling this blocking retire operation from the callback it is waiting for would deadlock.
Such APIs need deferred retirement or a nonblocking stop followed by later collection.
The Python model exhaustively enumerates short abstract traces of enter, leave, stop, and collect; it validates the state invariant, not a concrete mutex implementation or memory-ordering scheme.

\section{Attach the correct runtime on the correct thread}
\label{sec:ch73-threads}
A foreign-created thread is not automatically authorized to operate on managed objects.
Runtime attachment, interpreter identity, and synchronization are separate requirements.
In CPython, an attached thread state remains necessary on free-threaded builds, while GIL-enabled builds additionally serialize ordinary Python execution through the GIL~\cite{python-2026-capi}.
A convenience attachment API intended for the main interpreter must not be assumed to select an arbitrary subinterpreter.
Runtime shutdown creates another boundary: no callback may attempt attachment after its interpreter becomes unavailable.

Lua's state and API require an embedding policy that serializes access unless the embedder provides an appropriate alternative arrangement~\cite{ierusalimschy-2026-lua}.
The companion uses the calling thread and one active state; it does not dispatch callbacks from workers.
Its native pthread test instead calls independent scalar plans and never touches a managed runtime.
Table~\ref{tab:ch73-execution} distinguishes these verification claims.

\begin{table}[htbp]
\caption{Native parallelism and managed callback entry require different evidence.}
\label{tab:ch73-execution}
\centering
\begin{tabular}{ll}
\toprule
Path & Case-study contract \\
\midrule
Lua or Python callback & synchronous, attached calling thread \\
nested managed call & separate context; no shared callback slot \\
native pthread workers & independent plans; no runtime object access \\
asynchronous registration & protocol and finite model, not a native implementation \\
libffi closure fixture & live context until all direct test calls finish \\
\bottomrule
\end{tabular}
\end{table}

An event-queue design can move completion delivery back to the runtime's owning thread.
The worker then manipulates only native data and posts a retained completion record.
The queue still needs ownership transfer, shutdown draining, and cancellation rules; scheduling on the right thread solves only the entry-location part of the problem.

\section{Exercises}
\label{sec:ch73-exercises}
\begin{enumerate}
\item\label{ex:ch73-nested} Evaluate a visit of count three whose callback at $i$ visits count $i$ with callback $j+1$.
\item\label{ex:ch73-stop} Trace enter, stop, collect, leave, collect, treating invalid transitions as rejected requests that leave the state unchanged.
Which collection succeeds?
\item\label{ex:ch73-self} Explain the deadlock if an active callback synchronously waits for its own registration's active count to become zero.
\item\label{ex:ch73-thread} A worker holds a Python object pointer and a native mutex.
Why do those facts alone not permit a Python call, even on a free-threaded build?
\end{enumerate}

\section{Further reading}
\label{sec:ch73-reading}
libffi documents closure allocation and the lifetimes of closure metadata and user data~\cite{libffi-2026-manual}.
The Python thread-state documentation distinguishes interpreter attachment, the GIL, and free-threaded execution~\cite{python-2026-capi}.
Node-API's thread-safe function and asynchronous-work facilities illustrate explicit scheduling and ownership mechanisms for another runtime~\cite{node-2026-api}.
